
\documentclass[12pt]{ctexart}
\usepackage[a4paper,margin=1in]{geometry}
\usepackage{amsmath,amssymb,amsthm,mathtools,bm}
\usepackage{booktabs,longtable,array}
\usepackage{tabularx}
\usepackage{pdflscape}
\usepackage{enumitem}
\usepackage{setspace}
\usepackage[T1]{fontenc}
\usepackage[utf8]{inputenc}
\usepackage[hidelinks]{hyperref}

\newcolumntype{L}[1]{>{\raggedright\arraybackslash}p{#1}}
\newcolumntype{Y}{>{\raggedright\arraybackslash}X}

\setstretch{1.2}

\title{MAH Governance Cost Mechanism:\\
Model Mapping, Bellman Equations, and Empirical Implementation}
\author{}
\date{}

\begin{document}
\maketitle

\section{Purpose}

This note integrates two tasks into one coherent framework.

First, it gives a direct model mapping for the central mechanism in the MAH paper:
the reform lowers the governance and compliance friction of \emph{external commercialization}
for research-originated original-drug ideas.

Second, it updates the variable definitions and empirical implementation plan so that
the empirical section, the static model block, and the dynamic Bellman equations all
use the same object.

The key update is the following:

\[
\tau_{ext}
\]

is the central primitive through which MAH enters the model. It is \emph{not} a generic
policy dummy, and it is \emph{not} total governance cost for the whole firm. It is the
\emph{route-specific governance/compliance cost} of commercializing a research-originated
original-drug idea through an external qualified manufacturer while the innovator retains
authorization status.

\section{What \texorpdfstring{$\tau_{ext}$}{tau\_ext} is and what it is not}

\subsection{Definition}

Define
\[
\tau_{ext} \ge 0
\]
as the extra organizational, contractual, and compliance cost that a research-oriented
firm must bear when it keeps the project/authorization and completes commercialization
through an external qualified manufacturer rather than through full internal integration.

This cost includes:
\begin{itemize}[leftmargin=2em]
    \item drafting and enforcing entrustment agreements and quality agreements;
    \item auditing, monitoring, and coordinating the contract manufacturer;
    \item complying with holder-side quality responsibility and release obligations;
    \item managing cross-firm execution and accountability in the commercialization stage.
\end{itemize}

\subsection{What it does not include}

It does \emph{not} include:
\begin{itemize}[leftmargin=2em]
    \item physical manufacturing cost itself;
    \item R\&D cost of generating ideas;
    \item economy-wide demand shifts;
    \item a generic fixed operating cost of all firms.
\end{itemize}

Thus the correct interpretation is:

\[
\tau_{ext}
=
\text{governance/compliance cost of the external commercialization route}
\]

rather than

\[
\tau_{ext}
=
\text{all governance cost in the economy.}
\]

\section{Why this is the right MAH mechanism}

The economic role of MAH is not that it makes scientists more productive, and not that
it directly raises drug demand. Its distinctive role is organizational.

Before MAH, a research-originating firm without internal manufacturing capacity could
typically commercialize an original-drug idea only by integrating downstream production
or by relinquishing the project to an integrated producer. After MAH, the innovator may
retain authorization while contracting with a qualified external manufacturer, under a
legally recognized structure with explicit responsibility allocation and standardized
compliance requirements.

Therefore, the reform changes the cost of completing commercialization \emph{across firm
boundaries}. In model language, MAH lowers the route-specific governance/compliance
friction of external commercialization:

\[
\tau_{ext}^{post} < \tau_{ext}^{pre}.
\]

\section{Where \texorpdfstring{$\tau_{ext}$}{tau\_ext} enters the model}

\subsection{Core principle}

The key modeling principle is:

\[
\boxed{\tau_{ext}\text{ enters only at the stage where a research-originated idea is commercialized externally.}}
\]

It should not directly enter:
\begin{itemize}[leftmargin=2em]
    \item demand,
    \item scientific productivity,
    \item the physical production technology of manufacturers,
    \item generic fixed cost for all firms.
\end{itemize}

Instead, it enters the \emph{commercialization assignment problem} for original-drug ideas.

\section{Static block}

\subsection{External commercialization payoff for a type-B idea}

Let \(z\) denote firm capability and \(p_m\) the price of qualified contract manufacturing.

For a research-specialized firm \(B\), define the net payoff from commercializing one
original-drug idea through external qualified production as

\[
H_B^{out}(z,p_m;\tau_{ext})
=
\lambda_B(z,p_m)\,R_B(z)-p_m-\tau_{ext},
\]

where:
\begin{itemize}[leftmargin=2em]
    \item \(R_B(z)\) is the private value of a successfully commercialized original-drug idea;
    \item \(\lambda_B(z,p_m)\) is the success rate of external commercialization;
    \item \(p_m\) is the contract-manufacturing price;
    \item \(\tau_{ext}\) is the route-specific governance/compliance friction.
\end{itemize}

This separates two conceptually different costs:
\begin{align*}
p_m & : \text{market/production payment to the external manufacturer},\\
\tau_{ext} & : \text{organizational and legal cost of executing the route.}
\end{align*}

\subsection{Commercialization value per idea}

If the external route is the only route being emphasized, define
\[
\Psi_B(z;p_m,\tau_{ext})=\max\left\{H_B^{out}(z,p_m;\tau_{ext}),\,0\right\}.
\]

If an internal integration or project-transfer route is retained, then more generally:
\[
\Psi_B(z;p_m,\tau_{ext})
=
\max\left\{
H_B^{out}(z,p_m;\tau_{ext}),
H_B^{int}(z),
0
\right\}.
\]

The crucial point is that \(\tau_{ext}\) enters only the external route:
\[
\frac{\partial H_B^{out}}{\partial \tau_{ext}}=-1,
\qquad
\frac{\partial H_B^{int}}{\partial \tau_{ext}}=0.
\]

\subsection{Idea-generation problem of a type-B firm}

Let \(x_B \ge 0\) denote idea-generation intensity and let \(\eta>1\) govern convex idea-generation cost.
Then the optimized current-period payoff of a type-B firm is

\[
\Pi_B(z;p_m,\tau_{ext})
=
-f_B
+
\max_{x_B\ge 0}
\left\{
x_B\Psi_B(z;p_m,\tau_{ext})
-\frac{\chi_B}{\eta}x_B^\eta
\right\}.
\]

The first-order condition yields

\[
x_B^*(z;p_m,\tau_{ext})
=
\left(
\frac{\Psi_B(z;p_m,\tau_{ext})}{\chi_B}
\right)^{\frac{1}{\eta-1}}
\quad \text{if } \Psi_B>0.
\]

Hence the mechanism runs through:

\[
\tau_{ext}\downarrow
\Rightarrow
H_B^{out}\uparrow
\Rightarrow
\Psi_B\uparrow
\Rightarrow
x_B^*\uparrow.
\]

\section{Dynamic Bellman system}

\subsection{Type-B Bellman equation}

Let \(V_B(z)\) denote the value of operating as a type-B firm at state \(z\).
Suppose the firm can remain \(B\), switch to \(A\), switch to \(C\), or exit.

Define:
\[
W_{BB}(z)
=
\Pi_B(z;p_m,\tau_{ext})
+
\beta E\left[V_B(z')\mid z,B\right],
\]

\[
W_{BA}(z)
=
-\kappa_{BA}
+
\Pi_A(z;w)
+
\beta E\left[V_A(z')\mid z,A\right],
\]

\[
W_{BC}(z)
=
-\kappa_{BC}
+
\Pi_C(z;p_m,w)
+
\beta E\left[V_C(z')\mid z,C\right].
\]

Then:
\[
V_B(z)
=
\max\left\{
W_{BB}(z),\,
W_{BA}(z),\,
W_{BC}(z),\,
0
\right\}.
\]

The key point is:

\[
\tau_{ext} \text{ enters } V_B(z) \text{ through } W_{BB}(z) \text{ only.}
\]

That is, \(\tau_{ext}\) affects the value of remaining a research-specialized firm and using
the external commercialization route. It does not mechanically alter the value of becoming
integrated or becoming manufacturing-specialized.

\subsection{Type-A and Type-C Bellman equations}

For type-A:
\[
V_A(z)
=
\max\left\{
\begin{aligned}
&\Pi_A(z;w)+\beta E[V_A(z')\mid z,A],\\
&-\kappa_{AB}+\Pi_B(z;p_m,\tau_{ext})+\beta E[V_B(z')\mid z,B],\\
&-\kappa_{AC}+\Pi_C(z;p_m,w)+\beta E[V_C(z')\mid z,C],\\
&0
\end{aligned}
\right\}.
\]

For type-C:
\[
V_C(z)
=
\max\left\{
\begin{aligned}
&\Pi_C(z;p_m,w)+\beta E[V_C(z')\mid z,C],\\
&-\kappa_{CA}+\Pi_A(z;w)+\beta E[V_A(z')\mid z,A],\\
&-\kappa_{CB}+\Pi_B(z;p_m,\tau_{ext})+\beta E[V_B(z')\mid z,B],\\
&0
\end{aligned}
\right\}.
\]

Thus \(\tau_{ext}\) also affects \(A\) and \(C\) values \emph{indirectly}, through how it changes
the attractiveness of the \(B\) organizational state.

\subsection{Entry condition}

Let the value of entering as type-B be:
\[
V_B^{entry}
=
\int V_B(z_0)\,dG(z_0)-c_{EB}.
\]

Then lower \(\tau_{ext}\) raises type-B entry value:
\[
\tau_{ext}\downarrow
\Rightarrow
V_B(z)\uparrow
\Rightarrow
V_B^{entry}\uparrow.
\]

So the governance-cost mechanism affects both incumbent behavior and the entry margin.

\section{Observable variables and their corresponding model objects}

The table below updates the empirical implementation plan so that it is fully consistent
with the model object \(\tau_{ext}\).

\begingroup
\small
\setlength{\tabcolsep}{6pt}
\renewcommand{\arraystretch}{1.15}
\begin{landscape}
\begin{table}[htbp]
\centering
\caption{Observable variables and their mapping to model objects (Part I)}
\begin{tabularx}{\linewidth}{@{}L{0.16\linewidth}L{0.20\linewidth}YYY@{}}
\toprule
Variable & Mathematical definition & Economic meaning & Data source & Role in regression/calibration \\
\midrule
External commercialization adoption & \(\chi^{ext}_{ipt}=1\{h_{ipt}\neq m_{ipt}\}\) & Whether the product is commercialized through separated holder--manufacturer roles & NMPA product approvals, manufacturing licenses, provincial entrusted-manufacturing announcements & Directly observed route adoption; core mechanism variable \\
Type-B external commercialization & \(\chi^{B\to ext}_{ipt}\) & Most direct product-level realization of the MAH mechanism & Product-level holder--manufacturer matching plus firm classification & Sharpest empirical proxy for the external route \\
New original-drug products commercialized & \(N_t\) & Main policy outcome & CDE/NMPA application, approval, and launch data & Primary outcome variable \\
Type-B entry & \(Entry^B_t\) (count or share) & Upstream entry margin induced by lower external commercialization friction & Business registry, applicant/holder identity, firm classification & Auxiliary channel variable \\
Type-B conversion rate & \(Conv^B_t\) & Closest observable proxy for implementation success & CDE IND/project data + NMPA approvals + originator classification & Core proxy for identifying the effect of \(\tau_{ext}\) \\
\bottomrule
\end{tabularx}
\end{table}

\clearpage

{\footnotesize
\setlength{\tabcolsep}{4pt}
\begin{table}[htbp]
\centering
\caption{Observable variables and their mapping to model objects (Part II-A)}
\begin{tabularx}{\linewidth}{@{}L{0.15\linewidth}L{0.12\linewidth}YYY@{}}
\toprule
Variable & Mathematical definition & Economic meaning & Data source & Role in regression/calibration \\
\midrule
Pre-policy qualified manufacturing capacity & \(Cap_{r,d,0}\) & Exogenous supply-side shifter of external commercialization feasibility and governance friction & Manufacturing-license database, dosage/process qualification scope & Key heterogeneity variable; identifies where MAH should matter more \\
Producer-side outcomes & Revenue growth, gross margin, operating margin, entrusted-manufacturing revenue & How supply-side producers respond to higher external commercialization demand & Annual reports, prospectuses, firm databases & Supplementary equilibrium moments \\
Commercialization text proxies & Keyword frequency: entrusted manufacturing; MAH/holder; marketing release; CMO/CDMO & Supplementary evidence on organizational route adoption & Annual reports, prospectuses, announcements & Auxiliary evidence only; not a core variable \\
\bottomrule
\end{tabularx}
\end{table}

\begin{table}[htbp]
\centering
\caption{Observable variables and their mapping to model objects (Part II-B)}
\begin{tabularx}{\linewidth}{@{}L{0.15\linewidth}L{0.12\linewidth}YYY@{}}
\toprule
Variable & Mathematical definition & Economic meaning & Data source & Role in regression/calibration \\
\midrule
Generic cash-flow block & Generic sales, output, approvals & Controls for the role of generic production in firm cash flow and industry equilibrium & Annual reports, NMPA/CDE & Control and calibration block in two-medicine environment \\
Review-speed reform controls & Exposure to approval-delay reduction or drug-category-specific speed-up & Separates MAH from contemporaneous review reform & CDE/NMPA, drug-category data & Control block in reduced-form analysis \\
\bottomrule
\end{tabularx}
\end{table}
}
\end{landscape}
\endgroup

\section{What is identified in the data}

\subsection{Directly observed objects}

The model primitive \(\tau_{ext}\) is not directly observed.
What can be directly observed are product-level outcomes that are shaped by it.

The most important directly observed variables are:

\begin{enumerate}[leftmargin=2em]
    \item \textbf{External commercialization adoption}
    \[
    \chi^{ext}_{ipt}=1\{\text{holder}_{ipt}\neq \text{manufacturer}_{ipt}\}.
    \]

    \item \textbf{Type-B external commercialization}
    \[
    \chi^{B\to ext}_{ipt}=1\{\text{B-type holder uses an external qualified manufacturer}\}.
    \]

    \item \textbf{Type-B conversion rate}
    \[
    Conv^B_t
    =
    \frac{\#\{\text{B-originated original-drug products realized}\}}
         {\#\{\text{B-originated projects/applications}\}}.
    \]

    \item \textbf{New original-drug products commercialized}
    \[
    N_t = \#\{\text{new original-drug products commercialized at time } t\}.
    \]
\end{enumerate}

\subsection{Supply-side shifter for governance friction}

Construct a pre-policy measure of qualified manufacturing capacity:
\[
Cap_{r,d,0}
=
\#\{\text{qualified manufacturers in region } r \text{ and dosage/process class } d \text{ before MAH}\}.
\]

This variable is important because if MAH reduces the governance cost of external commercialization,
its effect should be larger where a suitable external manufacturing partner is easier to find and organize.

\section{Empirical implementation}

\subsection{Core reduced-form regressions}

\paragraph{Regression 1: route adoption}
\[
\chi^{B\to ext}_{ipt}
=
\alpha_i+\gamma_t+\beta Post_t+\delta X_{ipt}+\varepsilon_{ipt}.
\]

Expected sign:
\[
\beta>0.
\]

\paragraph{Regression 2: conversion rate}
\[
Conv^B_{rdt}
=
\alpha_{rd}+\gamma_t+\beta Post_t+\varepsilon_{rdt}.
\]

Expected sign:
\[
\beta>0.
\]

\paragraph{Regression 3: heterogeneity by external capacity}
\[
Y_{rdt}
=
\alpha_{rd}
+
\gamma_t
+
\beta_1 Post_t
+
\beta_2(Post_t\times Cap_{r,d,0})
+
\varepsilon_{rdt},
\]
where \(Y_{rdt}\) can be:
\begin{itemize}[leftmargin=2em]
    \item \(\chi^{B\to ext}\),
    \item \(Conv^B\),
    \item \(N_t\).
\end{itemize}

Expected sign:
\[
\beta_2>0.
\]

This is the sharpest reduced-form implication of the mechanism:
if MAH lowers the governance friction of external commercialization, then the reform should
have larger effects where the relevant external manufacturing route is easier to organize.

\subsection{Data implementation plan}

The preferred sequence is:

\begin{enumerate}[leftmargin=2em]
    \item Build a product-year dataset containing:
    \begin{itemize}[leftmargin=2em]
        \item product approval number,
        \item drug name, dosage form, and therapeutic category,
        \item original-drug indicator,
        \item holder/MAH identity,
        \item manufacturer identity,
        \item approval/marketing timing,
        \item province/city,
        \item manufacturer qualification scope.
    \end{itemize}

    \item Construct direct route-adoption variables:
    \[
    \chi^{ext}_{ipt}, \quad \chi^{B\to ext}_{ipt}.
    \]

    \item Construct implementation proxies:
    \[
    Conv^B_t, \quad N_t.
    \]

    \item Add the pre-policy supply-side shifter:
    \[
    Cap_{r,d,0}.
    \]

    \item Use firm annual reports and prospectuses only as supplementary validation for
    keywords such as:
    \begin{itemize}[leftmargin=2em]
        \item entrusted manufacturing,
        \item contract manufacturing,
        \item MAH/holder,
        \item product release,
        \item CMO/CDMO.
    \end{itemize}
\end{enumerate}

\section{How \texorpdfstring{$\tau_{ext}$}{tau\_ext} is calibrated}

\(\tau_{ext}\) should not be ``directly estimated'' from one variable.
Instead, it should be disciplined by a set of moments:

\begin{enumerate}[leftmargin=2em]
    \item external commercialization adoption rate;
    \item type-B external commercialization share;
    \item B-originated conversion rate;
    \item number of new original-drug products;
    \item heterogeneity with respect to \(Cap_{r,d,0}\);
    \item producer-side outcomes of qualified manufacturers.
\end{enumerate}

Thus the model treats \(\tau_{ext}\) as a latent structural friction that is pinned down by:
\begin{itemize}[leftmargin=2em]
    \item route adoption,
    \item route productivity,
    \item route capacity heterogeneity.
\end{itemize}

\section{Bottom line}

The updated structure can be summarized in one chain:

\[
MAH
\Rightarrow
\tau_{ext}\downarrow
\Rightarrow
H_B^{out}\uparrow
\Rightarrow
\Psi_B\uparrow
\Rightarrow
x_B\uparrow,\ V_B\uparrow,\ Entry_B\uparrow,\ N_t\uparrow.
\]

This has two advantages.

First, it gives MAH a single, sharp, and institutionally grounded entry point in the model.

Second, it aligns the theory with the empirical objects that can actually be measured:
route adoption, type-B conversion, new original-drug products, and heterogeneity by external
qualified manufacturing capacity.

\end{document}
